\setcounter{section}{6}
\section{平板中的穩態熱傳導}
\label{sec:ch7-slab}

\begin{quote}\small
本試作使用 S2：一份為工作流程測試而編造的兩頁英文來源。下文的「S2 第 201–202 頁」及方程、表格、題號都只屬於這份合成來源，不是任何真實教科書的頁碼或內容。
\end{quote}

若平板兩側的溫度已固定，怎麼由板厚與材料的熱傳導性質求出穿過平板的熱量傳遞速率？S2 第 201 頁把問題限制在一維、穩態的平板模型；先由溫度隨位置的變化得到熱通量，再乘上截面積，才得到整個截面的總熱傳率。讀完本節後，應能分辨熱通量與總熱傳率、由邊界溫度重建線性溫度分布，並用厚度改變解答 S2 的例題與習題。

\subsection{先定義模型與要算的量}
\sourcepage{S2（合成來源）p.201}

令座標 \(x\) 從熱面指向冷面，熱面在 \(x=0\)，冷面在 \(x=L\)。平板厚度為 \(L\)、固定截面積為 \(A\)；兩面的溫度分別是 \(T_H\) 與 \(T_C\)，且 \(T_H>T_C\)。位置 \(x\) 上的溫度記為 \(T(x)\)，邊界條件是 \(T(0)=T_H\) 和 \(T(L)=T_C\)。材料的熱導率 \(k\) 為常數且大於零。

這個模型採用一維穩態熱傳導：溫度只隨 \(x\) 改變，不隨時間改變；沒有內部熱生成，側面散熱可忽略。這些條件讓能量只需沿板的厚度方向追蹤。若有內部熱生成，或 \(k\) 隨位置改變，下面這段推導就不能原樣套用。

先分清兩個容易混為一談的量。熱通量（heat flux）\(q_x\) 是每單位面積、沿 \(+x\) 方向傳遞的熱量速率，單位為 \(\mathrm{W\,m^{-2}}\)；總熱傳率（total heat-transfer rate）\(\dot Q\) 是穿過整個截面的熱量速率，單位為 \(\mathrm W\)。在此固定面積且通量均勻的模型中，兩者由截面積相連；不能只因都叫「速率」就把它們當成同一量。

\begin{principle}{本節使用的模型條件}
一維、穩態；截面積 \(A\) 固定；\(k>0\) 且為常數；沒有內部熱生成；側面熱損失可忽略；邊界溫度固定為 \(T_H>T_C\)。熱通量的正方向定為從 \(x=0\) 指向 \(x=L\)。
\end{principle}

S2 將傅立葉定律（Fourier's law）作為經驗性的材料關係給定，而不是由能量守恆推導出來。它把局部熱通量連到溫度梯度：

\begin{equation}
q_x=-k\frac{dT}{dx}
\label{eq:7-1}
\end{equation}

負號的方向意義可直接用本題檢查。溫度由 \(x=0\) 的 \(T_H\) 降到 \(x=L\) 的 \(T_C\)，所以沿 \(+x\) 的溫度梯度為負；又因 \(k>0\)，式中的負號使 \(q_x>0\)。因此熱通量沿我們定義的 \(+x\) 方向，從高溫面指向低溫面。若漏掉負號，這個座標選擇下的方向就會相反。

\subsection{由能量守恆得到溫度分布}
穩態能量守恆、沒有內部熱生成且截面積固定時，S2 給出熱通量不隨 \(x\) 改變。再代入上面的傅立葉定律，並使用 \(k\) 為常數，溫度必須滿足：

\begin{align*}
\frac{dq_x}{dx}&=0,\\
\frac{d}{dx}\left(-k\frac{dT}{dx}\right)&=0,\\
\frac{d^2T}{dx^2}&=0 \qquad (k>0\text{ 為常數}).
\end{align*}

這裡的常數熱導率讓 \(k\) 能從對 \(x\) 的微分中提出；由於 \(k>0\)，要使熱通量梯度為零，就必須有溫度的二階導數為零。這並不是說溫度處處等於零，而是說其斜率不隨位置改變。因此溫度是 \(x\) 的一次函數。S2 直接給出結果，但略去了求兩個積分常數的代數步驟。

\begin{aisupplement}{由兩個邊界條件決定積分常數}
以下只補上 S2 第 201 頁省略的基本積分與代入，沒有另行推導傅立葉定律或能量守恆。從 \(d^2T/dx^2=0\) 積分一次，溫度梯度為常數 \(C_1\)；再積分一次，溫度為
\[
\frac{dT}{dx}=C_1,\qquad T(x)=C_1x+C_2.
\]
在熱面 \(x=0\)，邊界條件給 \(T_H=C_1(0)+C_2\)，故 \(C_2=T_H\)。在冷面 \(x=L\)，則
\[
T_C=C_1L+T_H
\quad\Longrightarrow\quad
C_1=\frac{T_C-T_H}{L}.
\]
這裡用到 \(L>0\)，才能除以厚度。兩個常數都已確定，將它們代回 \(T(x)=C_1x+C_2\)：
\[
T(x)=\left(\frac{T_C-T_H}{L}\right)x+T_H.
\]
把常數項 \(T_H\) 寫在前面，整理後與 S2 第 201 頁已給出的溫度分布一致：
\begin{equation}
T(x)=T_H+(T_C-T_H)\frac{x}{L}
\label{eq:7-2}
\end{equation}
\end{aisupplement}

式 \eqref{eq:7-2} 可用兩個端點直接核對：令 \(x=0\) 得 \(T_H\)，令 \(x=L\) 得 \(T_C\)。由於 \(T_C-T_H<0\)，斜率為負，故溫度從熱面到冷面線性下降。這個線性結果依賴剛才的模型條件；它不是不論內部是否產熱、熱導率是否改變都成立的通式。

\subsection{由溫度梯度到通量與總熱傳率}
現在可回到原問題：已知溫度分布，怎麼求穿過平板的熱傳？式 \eqref{eq:7-2} 的斜率在整塊板中固定，因此把它代入傅立葉定律，就得到每單位面積的熱通量。固定截面上的總熱傳率再由熱通量乘以面積：

\begin{equation}
\begin{aligned}
q_x&=-k\frac{T_C-T_H}{L}=\frac{k(T_H-T_C)}{L},\\
\dot Q&=Aq_x=\frac{kA(T_H-T_C)}{L}.
\end{aligned}
\label{eq:7-3}
\end{equation}

因 \(T_H-T_C>0\)，而 \(k,A,L\) 都為正，式 \eqref{eq:7-3} 的熱通量與總熱傳率皆為正，方向是 \(+x\)。兩者數值的差別來自所代表的物件：\(q_x\) 只按單位面積計，\(\dot Q\) 則把整個面積納入。量綱也不同：\(k(T_H-T_C)/L\) 的單位為 \(\mathrm{W\,m^{-2}}\)，再乘以 \(A\) 的 \(\mathrm{m^2}\) 才是 \(\mathrm W\)。

式 \eqref{eq:7-3} 也顯示各條件固定時的變化。若只把 \(L\) 加倍，熱通量與總熱傳率都減半；若增加 \(A\) 而維持同一溫度梯度，總熱傳率隨面積增加，熱通量則不變。若兩面溫度相同，式 \eqref{eq:7-2} 給出常數溫度，式 \eqref{eq:7-3} 給出零通量與零總熱傳率。這些結論都限於此處的穩態、無內部熱生成、常數 \(k\) 模型。

\subsection{例題 7.1：先求梯度，再分別算兩種速率}
\sourcepage{S2（合成來源）p.202，Example 7.1}

一塊平板的厚度 \(L=0.020\,\mathrm m\)，截面積 \(A=0.50\,\mathrm{m^2}\)，熱導率 \(k=0.80\,\mathrm{W\,m^{-1}K^{-1}}\)。熱面與冷面溫度分別為 \(30\,^{\circ}\mathrm C\) 與 \(20\,^{\circ}\mathrm C\)，且套用第 201 頁的模型。求：(a) 熱通量及其方向；(b) 總熱傳率。

\begin{examplebox}{例題 7.1 解答}
\textbf{判讀與選法。} 題目給了兩端溫度和位置，要求熱通量；傅立葉定律的輸入是溫度梯度，所以先用邊界溫度求斜率。求出通量後，再乘以整個截面積，才能回答總熱傳率。題目已明示採用本節模型，因此可使用式 \eqref{eq:7-2} 與 \eqref{eq:7-3}。

\textbf{計算 (a)。} 沿 \(+x\) 溫度由 \(30\,^{\circ}\mathrm C\) 降至 \(20\,^{\circ}\mathrm C\)，溫差為 \(10\,\mathrm K\)；兩個攝氏溫度相減時，溫差的數值與開爾文數值相同，故
\[
\frac{dT}{dx}=\frac{T_C-T_H}{L}
=\frac{(20-30)\,\mathrm K}{0.020\,\mathrm m}
=-500\,\mathrm{K\,m^{-1}}.
\]
將這個負斜率代入傅立葉定律，負負得正：
\[
q_x=-(0.80\,\mathrm{W\,m^{-1}K^{-1}})
(-500\,\mathrm{K\,m^{-1}})
=400\,\mathrm{W\,m^{-2}}.
\]
所以 (a) 的通量是 \(400\,\mathrm{W\,m^{-2}}\)，沿 \(+x\) 方向，亦即從熱面指向冷面。

\textbf{計算 (b)。} 總熱傳率包含整個 \(0.50\,\mathrm{m^2}\) 截面，故
\[
\dot Q=Aq_x
=(0.50\,\mathrm{m^2})(400\,\mathrm{W\,m^{-2}})
=200\,\mathrm W.
\]
因此 (b) 的總熱傳率為 \(200\,\mathrm W\)。單位相乘時 \(\mathrm{m^2}\) 抵銷 \(\mathrm{m^{-2}}\)，留下 \(\mathrm W\)；正值也與題設的高溫面在 \(x=0\)、熱流正方向指向冷面一致。這兩項分別回答了「每單位面積傳多少」與「整個截面傳多少」，不應互換。
\end{examplebox}

S2 的表 7.1 把同一例題的模型溫度列在三個位置。它是依式 \eqref{eq:7-2} 計算的值，不是實驗測量：

\begin{center}
\textbf{S2 表 7.1：例題 7.1 平板中的溫度}
\end{center}
\begin{center}
\begin{tabular}{rr}
\toprule
\(x\ /\ \mathrm m\) & \(T\ /\ ^{\circ}\mathrm C\) \\
\midrule
0.000 & 30 \\
0.010 & 25 \\
0.020 & 20 \\
\bottomrule
\end{tabular}
\end{center}

位置每增加 \(0.010\,\mathrm m\)，表中的溫度下降 \(5\,^{\circ}\mathrm C\)；中點 \(x=0.010\,\mathrm m\) 正好是兩端溫度的平均值。這些計算值呈現式 \eqref{eq:7-2} 的線性分布，支持的是本節模型內的結果，不能把它們稱為量測資料。

\begin{selftestbox}{習題 7.1（S2 Problem 7.1；提示與答案在章末）}
維持例題 7.1 的熱導率、截面積及兩面溫度不變，將平板厚度加倍為 \(0.040\,\mathrm m\)。求：(a) 新的熱通量；(b) 新的總熱傳率；(c) 新總熱傳率與例題 7.1 總熱傳率的比值。並解釋熱通量與總熱傳率為何有不同單位。
\end{selftestbox}
\sourcepage{S2（合成來源）p.202，Problem 7.1}
\subsection*{本章複習與題目對照}
本章的依賴順序是：指定模型與邊界溫度 \(\longrightarrow\) 穩態能量守恆使通量不隨 \(x\) 改變 \(\longrightarrow\) 得到線性溫度分布 \(\longrightarrow\) 由溫度梯度與傅立葉定律求熱通量 \(\longrightarrow\) 乘以截面積求總熱傳率。例題示範最後兩步的數值計算；習題改變板厚，練習沿同一條鏈重新計算並辨別量的單位。

\begin{center}
\begin{tabularx}{\linewidth}{p{0.20\linewidth} X X}
\toprule
學習目標 & 本章可回查處 & S2 題目對照 \\
\midrule
分辨每單位面積熱通量與整個截面的總熱傳率，並判讀正方向 & 定義段、式 \eqref{eq:7-1} 與 \eqref{eq:7-3} & Example 7.1 (a), (b)；Problem 7.1 (c) \\
由穩態條件與兩端溫度說明線性溫度分布 & 積分常數補步、式 \eqref{eq:7-2}、S2 表 7.1 & S2 §7.1 與 Example 7.1 \\
求熱通量、總熱傳率，並分析厚度加倍的影響 & 式 \eqref{eq:7-3} 及其固定其他量時的解讀 & Example 7.1 (a), (b)；Problem 7.1 (a)–(c) \\
\bottomrule
\end{tabularx}
\end{center}

\paragraph{帶條件的關係索引。}
\begin{itemize}
\item 傅立葉定律（式 \eqref{eq:7-1}）：S2 給定的經驗關係，將溫度梯度連到有方向的熱通量；此處 \(k>0\) 且固定。
\item 溫度分布（式 \eqref{eq:7-2}）：由無內部熱生成、穩態、常數 \(k\) 與兩個固定邊界溫度得到。
\item 通量與總熱傳率（式 \eqref{eq:7-3}）：在上述模型與固定截面下，\(q_x\) 按單位面積計，\(\dot Q=Aq_x\) 包含整個面積。
\end{itemize}

\paragraph{原書代表題與自測。}
S2 Example 7.1（p.202）示範由邊界溫度求梯度、再分別計算通量與總熱傳率，正文已完整解答。S2 Problem 7.1（p.202）改變厚度，測試三個相連問項；先自行作答，提示與完整參考解答位於本章末。這是 S2 提供的唯一習題，本試作不另加題目。

\paragraph{術語與符號。}
熱通量（heat flux）\(q_x\)：沿 \(+x\) 的單位面積熱傳率，\(\mathrm{W\,m^{-2}}\)；總熱傳率（total heat-transfer rate）\(\dot Q\)：整個截面的熱傳率，\(\mathrm W\)；熱導率（thermal conductivity）\(k\)：本模型中為正的常數，\(\mathrm{W\,m^{-1}K^{-1}}\)；\(T_H,T_C,T(x)\)：熱面、冷面及位置 \(x\) 的溫度；\(L\)：板厚，\(\mathrm m\)；\(A\)：截面積，\(\mathrm{m^2}\)。
\sourcepage{S2（合成來源）p.201–202}
